\documentclass[11pt]{article}

\usepackage{amsmath}
\usepackage{amsthm}
\usepackage{amssymb}
\usepackage{hyperref}

\newtheorem{definition}{Definition}
\newtheorem{lemma}{Lemma}
\newtheorem{theorem}{Theorem}
\newtheorem{corollary}{Corollary}
\newtheorem{proposition}{Proposition}
\theoremstyle{remark}
\newtheorem{remark}{Remark}
\newtheorem{example}{Example}

\title{Six Birds Theory}
\author{Ioannis Tsiokos}
\date{29 Jan 2026}

\begin{document}
\maketitle

\begin{abstract}
We present a math-only framework in which \emph{layers} arise as fixed points of idempotent operators and \emph{open-endedness} requires changing the operator rather than iterating it.
Order-theoretic closure operators and dynamics-induced endomaps $E_{\tau,f}$ (built from a Markov kernel, a coarse-graining lens, and a timescale) are treated uniformly via the abstraction of idempotent endomaps and their fixed points.
We define a computable total-variation idempotence defect for $E_{\tau,f}$ and show that small retention error implies approximate idempotence, yielding robust packaged “objects” at a given scale and language.
To prevent spurious directionality, we define arrow-of-time as path-space KL divergence and prove that coarse-graining cannot increase it (data processing); we further formalize the protocol-trap principle showing that autonomous protocol holonomy does not generate sustained asymmetry unless a genuine affinity is present in the lifted dynamics.
Finally, we give a finite forcing-style counting lemma: definable predicates are exponentially rare relative to a partition-language, so generic predicate extensions are overwhelmingly non-definable, providing a mathematically clean anti-saturation mechanism for strict ladder climbing.
\end{abstract}

\tableofcontents

% Scope IDs: see docs/spec/theorem-inventory.md and paper-contract.md
% Main-text claims include: T-CL-01, T-AOT-01, T-AOT-02, T-ACC-01, T-IC-01, T-IC-02, T-P1-01, T-P2-01, T-FOR-01, T-FOR-02.
\section{Introduction}
This paper develops a math-only framework for closure ladders and dynamics-induced endomaps on
finite state spaces. We connect order-theoretic closure to idempotent endomaps, quantify
path-space asymmetry via KL, and express autonomy/drive conditions as graph 1-form obstructions.
Computational evidence appears only in appendices.

\subsection{The organizing picture: a three-certificate loop}\label{sec:big-picture}

The results of this paper can be read as a minimal calculus for building hierarchical structure while keeping directionality honest.
The core observation is that three notions often conflated in informal discussions are logically distinct and have distinct mathematical \emph{certificates}:

\begin{itemize}
  \item \textbf{Stability / objects (packaging).} A layer is specified by an idempotent endomap (exact or approximate); its fixed points are the “objects” at that layer.
  In order theory this is a closure operator; in dynamics it is the induced endomap $E_{\tau,f}$.
  The certificate is (approximate) idempotence, quantified by the TV defect $\delta_{\tau,f}$.
  \item \textbf{Novelty / open-endedness (non-derivability).} Iterating a fixed closure/packaging rule saturates by idempotence.
  Strict ladder climbing therefore requires changing the language/closure rule.
  The certificate is \emph{non-definability}: a new predicate is not measurable with respect to the old partition-language.
  Our finite forcing lemma makes this generic: when there is hidden volume ($N-K>0$), uniformly random predicate extensions are overwhelmingly non-definable.
  \item \textbf{Directionality / irreversibility (audit).} Arrow-of-time is defined on path space as $\Sigma_T(\rho)=D_{\mathrm{KL}}(\mathbb P_{\rho,T}\|\mathcal R_*\mathbb P_{\rho,T})$.
  The certificate is either positive path-space KL or, under the accounted regime, a non-exact log-ratio 1-form (nonzero cycle affinities).
  Data processing ensures coarse-graining cannot create this certificate, and the protocol-trap theorem ensures it is not an artifact of hiding an internal clock.
\end{itemize}

These certificates compose into a simple loop:
choose a language (lens) $f$, induce a packaging endomap $E_{\tau,f}$ from dynamics and timescale, extract its (approximate) fixed points as objects, then adjoin a generic non-definable predicate to obtain a strict language extension and repeat—while auditing any apparent arrow-of-time via path-space KL and affinity.
The main text develops each piece of this loop in a representation-independent way and keeps all nonessential computational evidence in the appendices.

\paragraph{Roadmap.}
Section~\ref{sec:closure} develops the order-closure backbone and the saturation principle.
Section~\ref{sec:idempotent-endo} isolates idempotent endomaps as the unifying abstraction and introduces the induced empirical endomap $E_{\tau,f}$ with a computable idempotence defect and a qualified layer-birth criterion.
Section~\ref{sec:acc} formulates \textbf{A\_AUT}+\textbf{A\_REV}+\textbf{A\_ACC} as a graph 1-form exactness statement and defines affinities.
Section~\ref{sec:aot} defines arrow-of-time as path reversal asymmetry, proves data processing, and formalizes the protocol-trap/clock-audit principle (including the corrected autonomous reading of protocol holonomy).
Section~\ref{sec:forcing} gives the finite forcing lemma as an anti-saturation move for strict language extension.
Section~\ref{sec:primitives} summarizes the primitives P1--P6 as closure-changing operations, and Section~\ref{sec:examples} gives a few surgical examples.

\section{Related work}\label{sec:related}

This paper is written to be largely self-contained; we briefly situate four standard ingredients that
we use as mathematical tools.

\paragraph{Closure operators, reflections, and idempotents.}
Closure operators on preorders/lattices are classical objects in order theory: extensive, monotone,
and idempotent endomaps whose fixed points form a ``closed'' substructure.
In categorical language, (order-theoretic) closure operators are the thin-category shadow of
reflective subcategories and idempotent monads: a reflector is left adjoint to the inclusion of the
full subcategory of fixed points, and the universal property characterizes the ``least closed object
above'' a given one.
Our main text uses the order-theoretic formulation, and the Lean scaffold mechanizes the thin
(preorder) instance. See, e.g., \cite{DaveyPriestley2002,MacLane1998,Borceux1994}.

\paragraph{Coarse-graining of Markov dynamics and lumpability.}
State aggregation and lumpability for finite Markov chains provide exact hypotheses under which a
coarse observation of a Markov chain is itself Markov (or approximately so after subsampling).
We use coarse-graining primarily as a measurable map (a ``lens'') and emphasize statements that are
valid without assuming exact lumpability (notably, data-processing for relative entropy).
When we do speak about induced macro-dynamics, the relevant classical references are the standard
Markov chain texts that treat aggregation/lumpability and time reversal; see
\cite{KemenySnell1960,Norris1997}.

\paragraph{Path-space relative entropy and data processing.}
Relative entropy/KL divergence and the data processing inequality are standard in information theory:
pushforward under a measurable map cannot increase KL. We apply this at the level of path measures,
so that coarse observation cannot create time-reversal asymmetry.
The use of path-space relative entropy between forward and time-reversed Markov path measures as a
measure of irreversibility also appears in the literature on reversible Markov chains and (more
recently) stochastic thermodynamics; we cite these only as orientation since our arguments are finite
and measure-theoretic. See, e.g., \cite{CoverThomas2006,CsiszarKorner2011,Kelly1979,Seifert2012}.

\paragraph{Definability rarity and the forcing analogy.}
In our finite setting, ``definable from a language'' is treated as measurability with respect to a
finite $\sigma$-algebra (equivalently: being constant on the blocks of a partition). Our ``finite
forcing lemma'' is a counting statement: when blocks have nontrivial internal volume, predicates that
are definable from the coarse language form an exponentially small subset of all predicates, so a
random predicate extension is generically non-definable.
The term ``forcing'' is used only as an analogy to the role of generic extensions in set theory; our
results are finite and do not invoke infinitary semantics. Background references include standard
finite model theory and set-theory texts, e.g.\ \cite{Libkin2004,Cohen1963,Kunen2011}.

% Scope IDs: see docs/spec/theorem-inventory.md and paper-contract.md
\section{Conventions and notation}
\subsection{Finite state spaces, distributions, and kernels}

Throughout, $Z$ denotes a finite set (the microstate space). We write $\Delta(Z)$ for the probability simplex on $Z$.
We represent distributions as row vectors $\mu \in \Delta(Z)$ and Markov kernels as row-stochastic matrices $P$ so that the time-$1$ update is $\mu \mapsto \mu P$.
All logarithms are natural logs.

\subsection{Paths, time reversal, and relative entropy}

Fix a horizon $T \in \mathbb{N}$. A path is $z_{0:T} := (z_0,z_1,\dots,z_T)\in Z^{T+1}$.
Given an initial distribution $\rho\in\Delta(Z)$ and a time-homogeneous kernel $P$, the forward path law is
\[
\mathbb{P}_{\rho,T}(z_{0:T}) \;=\; \rho(z_0)\prod_{t=0}^{T-1} P(z_t,z_{t+1}).
\]
Let $R:Z^{T+1}\to Z^{T+1}$ be path reversal, $R(z_0,\dots,z_T)=(z_T,\dots,z_0)$, and let $R_*\mathbb{P}_{\rho,T}$ denote the pushforward law.

For probability measures $p,q$ on a finite set, we use the KL divergence
\[
D_{\mathrm{KL}}(p\|q) := \sum_x p(x)\log\frac{p(x)}{q(x)},
\]
with the conventions $0\log(0/q):=0$ and $p(x)>0,\,q(x)=0 \Rightarrow D_{\mathrm{KL}}(p\|q)=+\infty$.

\begin{definition}[Finite-horizon path reversal asymmetry]\label{def:path-asymmetry}
For any initial distribution $\rho$ and horizon $T$, define the path reversal asymmetry
\[
\Sigma_T(\rho) \;:=\; D_{\mathrm{KL}}\!\big(\mathbb{P}_{\rho,T}\,\|\,R_*\mathbb{P}_{\rho,T}\big)\in[0,\infty].
\]
\end{definition}

\begin{remark}[Stationarity is not needed for $\Sigma_T$ or data processing]
Definition~\ref{def:path-asymmetry} makes sense for any initial distribution $\rho$ (it may be $+\infty$ if reverse-support is missing).
The data processing inequality for KL is purely measure-theoretic and therefore applies to $\Sigma_T(\rho)$ without any stationarity assumption.
Stationarity enters only when one wants a steady-state interpretation (e.g.\ relating $\Sigma_T(\pi)$ to entropy production and/or considering long-time rates).
\end{remark}

\subsection{Coarse-graining lenses}

A (deterministic) coarse-graining or \emph{lens} is a map $f:Z\to X$ to a finite set $X$.
The pushforward of $\mu\in\Delta(Z)$ is $(f_*\mu)(x):=\sum_{z:f(z)=x}\mu(z)$.
On path space we use the coordinatewise lens $f^{T+1}:Z^{T+1}\to X^{T+1}$, and write $f_*\mathbb{P}_{\rho,T}$ for the induced macro path law.

\subsection{Support graphs and discrete 1-forms}

Given a kernel $P$ on $Z$, its directed support has an edge $z\to z'$ whenever $P(z,z')>0$.
When microreversibility holds (support is bidirected), we also use the induced undirected support graph with edge $\{z,z'\}$ whenever $P(z,z')>0$ and $P(z',z)>0$.

On bidirected support edges we use the antisymmetric edge 1-form
\[
a(z,z') \;:=\; \log\frac{P(z,z')}{P(z',z)}.
\]
For a directed cycle $\gamma=(z_0\!\to z_1\!\to\cdots\to z_m=z_0)$ in the bidirected support, its \emph{cycle integral} is
\[
\oint_\gamma a \;:=\; \sum_{i=0}^{m-1} a(z_i,z_{i+1}).
\]
If $a$ is exact, i.e.\ there exists $\Phi:Z\to\mathbb{R}$ with $a(z,z')=\Phi(z')-\Phi(z)$ on all bidirected edges, then $\oint_\gamma a=0$ for every cycle $\gamma$.

\subsection{Assumption bundles}

We will cite standing hypotheses by short tags (e.g.\ \textbf{A\_FIN}+\textbf{A\_LENS}) rather than repeating prose.
Use the following as the canonical assumption bundle list in this manuscript (tag spellings are aligned with \texttt{docs/spec/theorem-inventory.md}):

\begin{description}
  \item[\textbf{A\_FIN}] Finite state spaces (all measures and sums are finite; measurability is automatic).
  \item[\textbf{A\_AUT}] Autonomy: time-homogeneous dynamics; any protocol/phase variables are included in the state (no external schedule in main-text statements).
  \item[\textbf{A\_REV}] Microreversibility on support: $P(z,z')>0 \Rightarrow P(z',z)>0$ (bidirected support).
  \item[\textbf{A\_ACC}] Accounting/1-form decomposition: the antisymmetric log-ratio 1-form admits a decomposition into an exact part plus fixed affinity components (graph $H^1$ viewpoint).
  \item[\textbf{A\_NULL}] Null regime: all affinities vanish (the 1-form is exact / detailed-balance-like in the sense used here).
  \item[\textbf{A\_STAT}] A stationary distribution $\pi$ exists (full support when required by context).
  \item[\textbf{A\_LENS}] A specified coarse-graining lens $f:Z\to X$ is part of the statement.
\end{description}

% Scope IDs: see docs/spec/theorem-inventory.md and paper-contract.md
\section{Order-closure and closure ladders}\label{sec:closure}
\subsection{Order-theoretic closure and fixed points}

We begin with the minimal order-theoretic notion of closure. This section is intentionally
agnostic about dynamics, probability, or measurement; those enter later. The only idea we
need is that a closure operator defines a notion of ``completion,'' and its fixed points
constitute the ``objects'' at that level.

\begin{definition}[Closure operator]\label{def:closure-operator}
Let $(L,\le)$ be a preorder. A map $c:L\to L$ is a \emph{closure operator} if for all $x,y\in L$:
\begin{enumerate}
  \item (Extensive) $x \le c(x)$.
  \item (Monotone) $x \le y \Rightarrow c(x)\le c(y)$.
  \item (Idempotent) $c(c(x)) = c(x)$.
\end{enumerate}
\end{definition}

\begin{definition}[Closed points / objects]\label{def:closed-points}
Given a closure operator $c:L\to L$, its \emph{closed points} (or \emph{objects at level $c$}) are
the fixed points
\[
\mathrm{Fix}(c) \;:=\; \{x\in L : c(x)=x\}.
\]
\end{definition}

The order on $L$ induces a natural order on closure operators: we write $c \le d$ if
$c(x)\le d(x)$ for all $x\in L$ (pointwise order). Intuitively, $d$ is ``stronger'' than $c$ if it
closes more aggressively.

\begin{lemma}[Closed points are antitone in closure strength]\label{lem:closed-antitone}
Assume \textbf{A\_FIN}.
If $c \le d$ are closure operators on $L$, then $\mathrm{Fix}(d)\subseteq \mathrm{Fix}(c)$.
Equivalently: if $d(x)=x$ then $c(x)=x$.
\end{lemma}

\begin{proof}
Let $x\in \mathrm{Fix}(d)$, so $d(x)=x$. By extensiveness, $x\le c(x)$. By $c\le d$ we have
$c(x)\le d(x)=x$. Hence $x\le c(x)\le x$, so $c(x)=x$.
\end{proof}

\subsection{Closure ladders and saturation}

Idempotence implies that repeatedly applying a \emph{fixed} closure rule cannot yield unbounded novelty.

\begin{lemma}[One-step stabilization under iteration]\label{lem:closure-iterate-stabilizes}
Assume \textbf{A\_FIN}.
Let $c$ be a closure operator on $L$. Define the iterate $c^{(n)}$ by $c^{(0)}=\mathrm{id}$ and
$c^{(n+1)} = c\circ c^{(n)}$. Then for all $x\in L$ and all integers $n\ge 1$,
\[
c^{(n)}(x) = c(x).
\]
\end{lemma}

\begin{proof}
For $n=1$ the statement is trivial. If $c^{(n)}(x)=c(x)$ for some $n\ge 1$, then
$c^{(n+1)}(x)=c(c^{(n)}(x))=c(c(x))=c(x)$ by idempotence.
\end{proof}

\begin{corollary}[Closure saturates]\label{cor:closure-saturates}
Assume \textbf{A\_FIN}.
Fix a closure operator $c$. For any $x\in L$, the sequence
$x,\; c(x),\; c(c(x)),\; c(c(c(x))),\dots$ stabilizes after one closure step:
it is constant from the second term onward.
\end{corollary}

Thus, any notion of open-endedness based purely on repeated application of a single closure rule is
structurally limited. To obtain a strict hierarchy of ``new objects,'' one must change the closure rule.

\begin{definition}[Strictly stronger closure]\label{def:strictly-stronger}
For closure operators $c,d$ on $L$, write $c \prec d$ (``$d$ is strictly stronger than $c$'') if
$c\le d$ and there exists an $x\in L$ such that $d(x)\neq c(x)$.
\end{definition}

\begin{definition}[Closure ladder]\label{def:closure-ladder}
A \emph{closure ladder} is a sequence of closure operators $(c_n)_{n\in\mathbb N}$ on $L$ such that
$c_n \prec c_{n+1}$ for all $n$. The fixed-point sets form a nested chain
\[
\mathrm{Fix}(c_0)\supseteq \mathrm{Fix}(c_1)\supseteq \mathrm{Fix}(c_2)\supseteq \cdots
\]
by Lemma~\ref{lem:closed-antitone}.
\end{definition}

\begin{remark}[Lean formalization note]\label{rem:lean-closure}
The order-theoretic closure operator used here is formalized in the repository using mathlib's
\texttt{ClosureOperator}. The basic closure and ladder API lives in
\texttt{formal/ClosureLadder/Basic.lean}. In particular, the iterate-stabilization statement
Lemma~\ref{lem:closure-iterate-stabilizes} corresponds to Lean lemmas
\texttt{closure\*iterate\*succ} and \texttt{closure\*iterate\*ge\*one}, and ladder monotonicity
corresponds to \texttt{ladder\*mono}.
\end{remark}

\begin{remark}[Interpretation]
The key structural fact is that closure is idempotent: once a closure rule is fixed, repeated application cannot yield an unbounded novelty stream.
Mathematically, ``objects at a level'' are exactly fixed points, and a closure ladder can only climb by changing the closure rule (not by iterating the old one).
This is the precise content behind the slogan ``closure saturates.''
\end{remark}

% Scope IDs: see docs/spec/theorem-inventory.md and paper-contract.md
\section{Idempotent endomaps and dynamics-induced ``closures''}\label{sec:idempotent-endo}
\subsection{Idempotent endomaps}

The common mathematical nucleus behind both order-theoretic closures and the
dynamics-induced operators used later is \emph{idempotence}.  In this section we
package that nucleus in a way that avoids conflating two different notions:
\emph{order-closures} (which require an ambient preorder) versus \emph{idempotent
endomaps} (which require no order at all).

\begin{definition}[Idempotent endomap]
Let $A$ be a set (or type). An \emph{idempotent endomap} on $A$ is a function
$e : A \to A$ such that $e \circ e = e$, i.e.\ $e(e(a)) = e(a)$ for all $a\in A$.
\end{definition}

\begin{definition}[Fixed points]
For an endomap $e : A \to A$ define its fixed-point set
\[
\mathrm{Fix}(e) := \{a\in A : e(a)=a\}.
\]
We interpret $\mathrm{Fix}(e)$ as the collection of ``objects packaged by $e$''.
\end{definition}

\begin{lemma}[Idempotents split]
Assume \textbf{A\_FIN}.
Let $e : A \to A$ be an idempotent endomap. Let $i:\mathrm{Fix}(e)\hookrightarrow A$
denote inclusion, and define $r:A\to \mathrm{Fix}(e)$ by $r(a):=e(a)$ (viewed as an
element of $\mathrm{Fix}(e)$).
Then $r\circ i = \mathrm{id}_{\mathrm{Fix}(e)}$ and $i\circ r = e$.
In particular, $e$ restricts to the identity on $\mathrm{Fix}(e)$ and
$\mathrm{Fix}(e)$ canonically identifies with the image of $e$.
\end{lemma}

\begin{proof}
For $a\in A$, idempotence gives $e(e(a))=e(a)$, so $e(a)\in \mathrm{Fix}(e)$ and
$r$ is well-defined. For $x\in \mathrm{Fix}(e)$ we have $r(i(x))=e(x)=x$, hence
$r\circ i=\mathrm{id}$. Finally, for $a\in A$, $(i\circ r)(a)=i(e(a))=e(a)$ by
definition of $i$, so $i\circ r=e$.
\end{proof}

\begin{remark}[Order-closures as a special case]
If $L$ is a preorder and $c:L\to L$ is an order-closure (extensive, monotone,
idempotent), then forgetting the order yields an idempotent endomap on the
underlying set. The converse is false: an idempotent endomap need not be monotone
or extensive with respect to any nontrivial order.  In particular, the
dynamics-induced operators $E_{\tau,\Pi}$ (defined later on probability simplices)
are treated as \emph{(approximate) idempotent endomaps} with fixed points and
idempotence defects measured in a chosen metric, rather than as order-closures.
\end{remark}

\begin{remark}[Formalization note]
The constructions in this subsection are formalized in Lean.
See \texttt{formal/ClosureLadder/IdempotentEndo.lean} for the definitions
\texttt{ClosureLadder.IdempotentEndo} and \texttt{ClosureLadder.IdempotentEndo.FixedPoints},
and the bridge lemma \texttt{ClosureLadder.ClosureOperator.toIdempotentEndo}.
\end{remark}

\subsection{Dynamics-induced empirical endomaps}
\label{sec:empirical-closure}

Throughout this section we assume \textbf{A\_FIN} and fix a finite state space $Z$.
Let $P$ be a (row-stochastic) Markov kernel on $Z$, and let $\Delta(Z)$ denote the simplex
of probability distributions on $Z$ (row vectors).
Fix a \emph{lens} $f:Z\to X$ (assumption \textbf{A\_LENS}), with fibers (blocks)
$B_x:=f^{-1}(x)$.

\begin{definition}[Coarse map and canonical lift; \textbf{D-IC-01}]
\label{def:D-IC-01}
Let $Q_f:\Delta(Z)\to \Delta(X)$ be the pushforward (coarse-graining) map
\[
(Q_f\mu)(x) \;:=\; \sum_{z\in B_x}\mu(z).
\]
Fix a \emph{canonical lift} $U_f:\Delta(X)\to \Delta(Z)$ specified by a choice of
\emph{prototype} distributions $(u_x)_{x\in X}$ with $\mathrm{supp}(u_x)\subseteq B_x$ and
$\sum_{z\in B_x}u_x(z)=1$, by setting
\[
(U_f\nu)(z) \;:=\; \sum_{x\in X}\nu(x)\,u_x(z).
\]
Equivalently, $U_f\delta_x=u_x$ and $Q_f U_f = \mathrm{id}_{\Delta(X)}$.

For an integer time scale $\tau\ge 1$, define the \emph{induced empirical endomap}
\[
E_{\tau,f}:\Delta(Z)\to \Delta(Z),
\qquad
E_{\tau,f}(\mu) \;:=\; U_f\bigl(Q_f(\mu P^\tau)\bigr).
\]
\end{definition}

\begin{remark}[Not an order-closure]
The map $E_{\tau,f}$ is an affine endomap of $\Delta(Z)$ (indeed a Markov kernel on $Z$),
and it is always idempotent as an \emph{endomap} only in special regimes.
We treat $E_{\tau,f}$ within the ``fixed points of (approx) idempotents'' framework
(\S\ref{sec:idempotent-endo}), not as an order-closure operator.
See the repo spec note \texttt{docs/spec/empirical-vs-order-closure.md}.
\end{remark}

\begin{definition}[TV idempotence defect; \textbf{D-IC-02}]
\label{def:D-IC-02}
Let $\|\cdot\|_{\mathrm{TV}}$ denote total variation on $\Delta(Z)$, i.e.
$\|\mu-\nu\|_{\mathrm{TV}}=\tfrac12\|\mu-\nu\|_1$.
The \emph{idempotence defect} of $E_{\tau,f}$ is
\[
\delta_{\tau,f}
\;:=\;
\sup_{\mu\in \Delta(Z)} \bigl\|E_{\tau,f}(E_{\tau,f}(\mu)) - E_{\tau,f}(\mu)\bigr\|_{\mathrm{TV}}.
\]
Since $\mu\mapsto \| \mu M\|_1$ is convex for fixed $M$ and $\Delta(Z)$ is the convex hull
of the Dirac masses $(\delta_z)_{z\in Z}$, the supremum is attained on extreme points:
\[
\delta_{\tau,f}
\;=\;
\max_{z\in Z} \bigl\|\delta_z(E_{\tau,f}^2 - E_{\tau,f})\bigr\|_{\mathrm{TV}}
\;=\;
\frac12\max_{z\in Z}\sum_{z'\in Z}\bigl|(E_{\tau,f}^2 - E_{\tau,f})(z,z')\bigr|.
\]
\end{definition}

\begin{definition}[Prototype stability at scale $\tau$]
For each $x\in X$ define the (scale-$\tau$) stability of the prototype $u_x$ by
\[
s_{\tau,f}(x) \;:=\; \bigl\|E_{\tau,f}(u_x)-u_x\bigr\|_{\mathrm{TV}}
\;=\;
\bigl\|U_f(Q_f(u_xP^\tau)) - U_f(\delta_x)\bigr\|_{\mathrm{TV}}.
\]
We say $u_x$ is \emph{$\varepsilon$-stable} if $s_{\tau,f}(x)\le \varepsilon$.
\end{definition}

\begin{theorem}[Approximate idempotence from prototype retention; \textbf{T-IC-02}]
\label{thm:T-IC-02}
Assume \textbf{A\_FIN}+\textbf{A\_LENS} and fix $E_{\tau,f}$ as in Definition~\ref{def:D-IC-01}.
Define a \emph{retention error}
\[
\varepsilon_{\tau,f}
\;:=\;
\max_{x\in X}\bigl\|Q_f(u_xP^\tau)-\delta_x\bigr\|_{\mathrm{TV}}.
\]
Then:
\begin{enumerate}
\item (Prototype stability bound) For all $x\in X$, $s_{\tau,f}(x)\le \varepsilon_{\tau,f}$.
\item (Idempotence defect bound) $\delta_{\tau,f}\le \varepsilon_{\tau,f}$.
In particular, if $\varepsilon_{\tau,f}=0$ then $E_{\tau,f}$ is exactly idempotent.
\end{enumerate}
\end{theorem}

\begin{proof}
Since $U_f$ is a stochastic map, it is a contraction in total variation.
Thus for each $x$,
\[
\|E_{\tau,f}(u_x)-u_x\|_{\mathrm{TV}}
=
\|U_f(Q_f(u_xP^\tau)) - U_f(\delta_x)\|_{\mathrm{TV}}
\le
\|Q_f(u_xP^\tau)-\delta_x\|_{\mathrm{TV}}
\le \varepsilon_{\tau,f}.
\]
This proves (1).

For (2), write any $\mu\in\Delta(Z)$ as $\mu_1:=E_{\tau,f}(\mu)=\sum_{x\in X}w_x u_x$
where $w:=Q_f(\mu P^\tau)\in \Delta(X)$ (so $w_x\ge 0$ and $\sum_x w_x=1$).
Then by linearity,
\[
E_{\tau,f}(\mu_1) - \mu_1
=
\sum_{x\in X} w_x\bigl(E_{\tau,f}(u_x)-u_x\bigr).
\]
Taking TV norms and using the bound from (1),
\[
\|E_{\tau,f}(\mu_1) - \mu_1\|_{\mathrm{TV}}
\le
\sum_{x\in X}w_x \|E_{\tau,f}(u_x)-u_x\|_{\mathrm{TV}}
\le
\sum_{x\in X}w_x \varepsilon_{\tau,f}
=
\varepsilon_{\tau,f}.
\]
Now take the supremum over $\mu$ to obtain $\delta_{\tau,f}\le \varepsilon_{\tau,f}$.
\end{proof}

\begin{theorem}[Refinement can help or hurt; \textbf{T-IC-01}]
\label{thm:T-IC-01}
Assume \textbf{A\_FIN}+\textbf{A\_LENS}.
There exist finite Markov kernels $P$ and times $\tau$ together with two lenses
$f_{\mathrm{fine}}:Z\to X_{\mathrm{fine}}$ and $f_{\mathrm{coarse}}:Z\to X_{\mathrm{coarse}}$
such that $f_{\mathrm{coarse}}$ is a coarsening of $f_{\mathrm{fine}}$ (the fine partition refines the coarse),
and for a fixed stability threshold $\varepsilon>0$ the number of $\varepsilon$-stable prototypes
for $E_{\tau,f}$ can either increase or decrease under refinement.

More precisely, both behaviors occur:
\begin{enumerate}
\item (\emph{Refinement reveals objects}) in one example, refinement increases the number of
$\varepsilon$-stable prototypes.
\item (\emph{Refinement destroys objects}) in another example, refinement decreases the number of
$\varepsilon$-stable prototypes.
\end{enumerate}
\end{theorem}

\begin{proof}[Proof idea by explicit two-block constructions]
The mechanism is controlled by the retention errors $\varepsilon_{\tau,f}$ of
Theorem~\ref{thm:T-IC-02}.

Consider a state space split into two blocks $B_0,B_1$ and a lift $U_f$ whose prototypes
$u_0,u_1$ are supported on the blocks. If the dynamics jump from $B_0$ to $B_1$ with
probability $p$ each step (and otherwise remain in the current block), then starting from $u_0$,
the probability of exiting $B_0$ within $\tau$ steps is $1-(1-p)^\tau$.
Hence $\|Q_f(u_0P^\tau)-\delta_0\|_{\mathrm{TV}}$ is on the order of $1-(1-p)^\tau$,
and similarly for $u_1$.

Now take a coarse lens that merges two metastable regions into one block, and a refinement that
splits them. If $p$ is very small (relative to $\tau$), the refined prototypes have small
retention error and are stable, so refinement increases the stable count.
If $p$ is not small (relative to $\tau$), the refined prototypes rapidly mix across the split
boundary, their retention errors exceed the stability threshold, and refinement decreases the
stable count.  \qedhere
\end{proof}

\begin{remark}[Qualified ``layer birth'']
The induced endomap $E_{\tau,f}$ packages a distribution by evolving for $\tau$ steps and then
forgetting micro-detail below the lens $f$, re-instantiating it via the chosen prototypes.
Theorem~\ref{thm:T-IC-02} shows that a small retention error implies small idempotence defect.
In this regime, the (approximately) stable prototypes $u_x$ behave as emergent ``objects'' at
scale $\tau$ in the language $f$.
Refinement changes the lens and can therefore reveal or destroy such objects
(Theorem~\ref{thm:T-IC-01}).
\end{remark}

\begin{remark}[Interpretation]
Idempotent endomaps isolate the common nucleus behind both order-closures and dynamics-induced packaging: fixed points are the packaged objects, and idempotence means reapplying the packaging does nothing new.
For empirical operators $E_{\tau,f}$ the right question is therefore not ``is this a closure operator?'' but ``how close is it to idempotent, and what are its stable fixed points at the chosen scale/lens?''
Choosing a computable defect (TV/L1) turns this into an exact diagnostic rather than a metaphor.
\end{remark}

% Scope IDs: see docs/spec/theorem-inventory.md and paper-contract.md
\section{AUT + REV + ACC regime and graph 1-forms}
\label{sec:acc}

This section isolates the representation-independent content behind the admissibility regime
\textbf{A\_AUT}+\textbf{A\_REV}+\textbf{A\_ACC} that we use throughout: the obstruction to
``no-drive'' behavior is a cohomological (cycle) obstruction on the bidirected support graph.
Our presentation is discrete-time and purely graph-theoretic; any continuous-time (CTMC) reading
is optional and out of scope for the main text.

\subsection{Bidirected support and the log-ratio 1-form}

Assume \textbf{A\_FIN} and let $P$ be a time-homogeneous Markov kernel on a finite state space $Z$.
Write $P_{ij}:=P(i,j)$.
Under \textbf{A\_REV} (microreversibility on support), we have
\[
P_{ij}>0 \implies P_{ji}>0,
\]
so every allowed transition comes with a reverse transition.

Define the (undirected) \emph{support graph} $G=(Z,E)$ by declaring $\{i,j\}\in E$ iff $P_{ij}>0$
(equivalently $P_{ji}>0$ by \textbf{A\_REV}). Let $\vec E$ be the corresponding set of
\emph{oriented} edges: for each $\{i,j\}\in E$ we include both $(i,j)$ and $(j,i)$.

\begin{definition}[Edge log-ratio 1-form]
\label{def:edge-one-form}
Assume \textbf{A\_FIN}+\textbf{A\_REV}. The \emph{antisymmetric edge 1-form} associated to $P$ is
the function $a:\vec E\to \mathbb R$ defined by
\[
a_{ij} := a(i,j) := \log\frac{P_{ij}}{P_{ji}} \qquad \text{for } (i,j)\in \vec E.
\]
Then $a_{ji}=-a_{ij}$.
\end{definition}

\begin{remark}[Why \textbf{A\_REV} matters]
If \textbf{A\_REV} fails, then ratios such as $P_{ij}/P_{ji}$ may be undefined, and many
time-reversal quantities become infinite rather than merely nonzero. In this paper we treat
\textbf{A\_REV} as the basic finiteness/admissibility condition for edge log-ratios.
\end{remark}

\subsection{Cycle integrals, exactness, and the null regime}

A (simple) cycle in $G$ is a sequence $(v_0,v_1,\dots,v_m=v_0)$ such that $\{v_k,v_{k+1}\}\in E$
for all $k$.

\begin{definition}[Cycle integral / affinity along a cycle]
\label{def:cycle-integral}
Given an antisymmetric edge 1-form $a$ and an oriented cycle $\gamma=(v_0,\dots,v_m=v_0)$, define
\[
\mathcal A(\gamma) := \sum_{k=0}^{m-1} a_{v_k v_{k+1}}.
\]
Reversing the orientation negates $\mathcal A(\gamma)$.
\end{definition}

\begin{definition}[Exact 1-forms]
\label{def:exact-one-form}
An antisymmetric 1-form $a$ on $\vec E$ is \emph{exact} if there exists a potential
$\Phi:Z\to\mathbb R$ such that for every $(i,j)\in \vec E$,
\[
a_{ij} = \Phi(j)-\Phi(i).
\]
We write $a=d\Phi$ in this case.
\end{definition}

\begin{lemma}[Exact forms have zero cycle integrals]
\label{lem:exact-zero-cycles}
Assume \textbf{A\_FIN}+\textbf{A\_REV}.
If $a=d\Phi$ is exact, then $\mathcal A(\gamma)=0$ for every cycle $\gamma$ in $G$.
\end{lemma}

\begin{proof}
Along a cycle $\gamma=(v_0,\dots,v_m=v_0)$ we have
\[
\mathcal A(\gamma)=\sum_{k=0}^{m-1}\big(\Phi(v_{k+1})-\Phi(v_k)\big)=\Phi(v_m)-\Phi(v_0)=0.
\]
\end{proof}

\begin{theorem}[Cycle criterion for exactness]
\label{thm:cycle-criterion-exact}
Assume \textbf{A\_FIN}+\textbf{A\_REV} and let $a$ be an antisymmetric 1-form on the oriented edge
set $\vec E$ of $G$. The following are equivalent:
\begin{enumerate}
\item $a$ is exact: $a=d\Phi$ for some $\Phi:Z\to\mathbb R$.
\item $\mathcal A(\gamma)=0$ for every cycle $\gamma$ in $G$.
\item $\mathcal A(\gamma)=0$ for every cycle $\gamma$ in some fixed cycle basis of $G$
(e.g.\ a fundamental cycle basis from a spanning tree).
\end{enumerate}
\end{theorem}

\begin{proof}[Proof sketch]
(1)$\Rightarrow$(2) is Lemma~\ref{lem:exact-zero-cycles}. (2)$\Rightarrow$(3) is immediate.

For (3)$\Rightarrow$(1), fix a spanning tree $T$ of each connected component of $G$.
Choose a root $r$ in each component and define $\Phi(r)=0$; for any vertex $v$ in that component,
let $\Phi(v)$ be the sum of $a$ along the unique tree path from $r$ to $v$ (with signs determined
by orientations).
This makes $a_{ij}=\Phi(j)-\Phi(i)$ hold for every tree edge $(i,j)$.
For any chord edge $e=(u,v)\notin T$, the fundamental cycle $\gamma_e$ is the tree path from $u$
to $v$ together with $e$; the hypothesis $\mathcal A(\gamma_e)=0$ forces $a_{uv}=\Phi(v)-\Phi(u)$
on the chord as well.
Thus $a=d\Phi$ on all of $\vec E$.
\end{proof}

\begin{corollary}[Null regime as ``no-drive'' / detailed-balance-like]
\label{cor:null-regime}
In the \textbf{A\_NULL} regime (all cycle integrals vanish), the log-ratio 1-form $a$ is exact.
Equivalently, there exists $\Phi:Z\to\mathbb R$ such that
\[
\log\frac{P_{ij}}{P_{ji}} = \Phi(j)-\Phi(i) \quad \text{for all } (i,j)\in \vec E.
\]
In particular, defining $\pi(i)\propto e^{\Phi(i)}$ yields the detailed balance relations
$\pi(i)P_{ij}=\pi(j)P_{ji}$ on support.
\end{corollary}

\subsection{Accounting as coordinates on cycle space}

Let $\beta_1(G)$ be the cycle rank of the undirected support graph.
Choosing a cycle basis $(\gamma_1,\dots,\gamma_{\beta_1})$ defines a coordinate vector of
\emph{affinities}
\[
A_r := \mathcal A(\gamma_r), \qquad r=1,\dots,\beta_1(G).
\]
Changing the cycle basis changes coordinates but not the underlying obstruction: $A=0$ in one
basis iff $A=0$ in every basis (Theorem~\ref{thm:cycle-criterion-exact}).

We regard \textbf{A\_ACC} as the decision to work with this decomposition viewpoint:
the antisymmetric log-ratio 1-form $a$ splits into an exact (gauge) part plus a cycle/affinity
part, and the affinity part vanishes exactly in the null regime \textbf{A\_NULL}.

\begin{remark}[Implementation note for Appendix C]
The repository includes a small graph-theoretic routine that computes cycle affinities and tests
exactness on finite Markov kernels (reversible vs biased 3-cycle examples). We relegate all such
numerical evidence to Appendix~C per the paper contract.
\end{remark}

\begin{remark}[Interpretation]
The accounting regime is, at its core, a statement about a discrete 1-form on the bidirected support graph.
``Drive'' is not a vague physical intuition here: it is exactly the cohomological obstruction to exactness, detected by nonzero cycle integrals.
Affinities are simply coordinates for that obstruction in a chosen cycle basis; the null regime is the exact (no-cycle-obstruction) case.
\end{remark}

% Scope IDs: see docs/spec/theorem-inventory.md and paper-contract.md
\section{Arrow-of-time as path reversal asymmetry; data processing; protocol trap}\label{sec:aot}

Throughout this section we assume \textbf{A\_FIN}. Let $Z$ be a finite state space and let
$P:Z\times Z\to [0,1]$ be a time-homogeneous Markov kernel.

\begin{definition}[Forward path law and reversal (D-AOT-01)]
Fix a horizon $T\in \mathbb N$ with $T\ge 1$ and an initial distribution $\rho\in\Delta(Z)$.
The (forward) length-$T$ path law $\mathbb P_{\rho,T}$ is the probability measure on $Z^{T+1}$ defined by
\[
\mathbb P_{\rho,T}(z_0,\dots,z_T)
:= \rho(z_0)\prod_{t=0}^{T-1} P(z_t,z_{t+1}).
\]
Let $\mathcal R:Z^{T+1}\to Z^{T+1}$ be the reversal map $\mathcal R(z_0,\dots,z_T)=(z_T,\dots,z_0)$, and write $\mathcal R_*\mu$ for pushforward of a measure $\mu$.
Define the (finite-horizon) \emph{path reversal asymmetry}
\[
\Sigma_T(\rho)\;:=\; D_{\mathrm{KL}}\!\left(\mathbb P_{\rho,T}\,\middle\|\,\mathcal R_*\mathbb P_{\rho,T}\right)\in[0,\infty].
\]
We use the standard convention $D_{\mathrm{KL}}(p\|q)=+\infty$ if $p$ is not absolutely continuous with respect to $q$.
\end{definition}

\begin{remark}[Stationarity is not required (L-AOT-STAT-01)]
The quantity $\Sigma_T(\rho)$ is defined for any initial distribution $\rho$; it can be strictly positive even in a detailed-balance chain when $\rho$ is not stationary (a boundary-term effect).
Stationarity is only needed when we want a \emph{steady-state} interpretation (entropy production rate), not for the definition of $\Sigma_T$ or for the data-processing inequality below.
\end{remark}

\begin{definition}[Steady-state entropy production (D-AOT-EP-01)]
Assume additionally \textbf{A\_STAT} and let $\pi$ be a stationary distribution for $P$.
The (finite-horizon) steady-state entropy production is $\mathrm{EP}_T:=\Sigma_T(\pi)$.
When a per-step rate is desired, we write $\mathrm{ep}:=\tfrac1T\Sigma_T(\pi)$ (and, under standard ergodicity conditions, one can pass to the limit $T\to\infty$).
\end{definition}

\subsection{Data processing: coarse-graining cannot create irreversibility}

Now assume \textbf{A\_LENS} and fix a coarse-graining (lens) $f:Z\to X$.
Let $f^{T+1}:Z^{T+1}\to X^{T+1}$ be the coordinatewise map $f^{T+1}(z_0,\dots,z_T)=(f(z_0),\dots,f(z_T))$.
Define the observed path law $\mathbb Q_{\rho,T}:=(f^{T+1})_*\mathbb P_{\rho,T}$.

\begin{theorem}[Data processing for path reversal asymmetry (T-AOT-01)]\label{thm:dpi_path}
Under \textbf{A\_FIN}+\textbf{A\_LENS}, for every $T\ge 1$ and every initial distribution $\rho$,
\[
D_{\mathrm{KL}}\!\left(\mathbb Q_{\rho,T}\,\middle\|\,\mathcal R_*\mathbb Q_{\rho,T}\right)
\;\le\;
D_{\mathrm{KL}}\!\left(\mathbb P_{\rho,T}\,\middle\|\,\mathcal R_*\mathbb P_{\rho,T}\right)
=\Sigma_T(\rho).
\]
\end{theorem}

\begin{proof}[Proof sketch]
The reversal map $\mathcal R$ commutes with coordinatewise coarse-graining: $\mathcal R\circ f^{T+1}=f^{T+1}\circ \mathcal R$.
Hence $\mathcal R_*\mathbb Q_{\rho,T}=(f^{T+1})_*(\mathcal R_*\mathbb P_{\rho,T})$.
The claim is then the usual contraction of KL under measurable maps: for any measurable $g$, $D_{\mathrm{KL}}(g_*p\|g_*q)\le D_{\mathrm{KL}}(p\|q)$.
\end{proof}

\begin{corollary}[No false positives from forgetting variables]
Assume \textbf{A\_FIN}+\textbf{A\_LENS}.
If the lifted system has zero path reversal asymmetry at horizon $T$, then every coarse observation has zero path reversal asymmetry at the same horizon.
Coarse-graining can hide irreversibility (false negatives), but it cannot create it.
\end{corollary}

\subsection{Protocol trap: apparent stroboscopic arrows and the ``clock audit''}

We now pin down a class of autonomous protocol models and the sense in which ``P3 requires P6''.
Let $X$ be a microstate space and let $\Phi=\{0,1,\dots,m-1\}$ be a finite phase space.
A \emph{protocol family} is a collection of kernels $\{K_\phi\}_{\phi\in\Phi}$ on $X$ together with a phase kernel $S$ on $\Phi$.

\begin{definition}[Autonomous $m$-phase random-scan protocol (D-PROT-02)]
Define an autonomous Markov chain on the product space $Z:=X\times \Phi$ by
\[
P\big((x,\phi),(x',\phi')\big)\;:=\; S(\phi,\phi')\,K_\phi(x,x').
\]
We call this the \emph{lifted} (phase-included) protocol model.
\end{definition}

\begin{theorem}[Protocol trap / hidden-drive principle (T-AOT-02)]\label{thm:protocol_trap}
Assume \textbf{A\_FIN}+\textbf{A\_AUT}.
Suppose:
(i) the phase chain $S$ is reversible with stationary distribution $s$;
(ii) each $K_\phi$ is reversible with respect to a \emph{common} stationary distribution $\pi$ on $X$.
Then the product measure $\mu:=\pi\times s$ is stationary for the lifted chain on $Z=X\times\Phi$ and the lifted chain is reversible (hence $\Sigma_T(\mu)=0$ for all $T$).
Consequently, by Theorem~\ref{thm:dpi_path} applied to the projection $X\times\Phi\to X$, the observed process on $X$ has zero path reversal asymmetry at every finite horizon.

In particular, any \emph{apparent} arrow-of-time obtained by analysing a stroboscopic composition of kernels on $X$ must be traced to either:
(a) an external schedule (violating \textbf{A\_AUT}), or
(b) a driven/bias phase dynamics (nonzero affinity, cf.\ Section~\ref{sec:acc}).
\end{theorem}

\begin{corollary}[``P3 needs P6'' under autonomy (C-ACC-01)]
Under \textbf{A\_AUT}+\textbf{A\_ACC}+\textbf{A\_NULL} (no affinities anywhere in the lifted state space), protocol holonomy alone does not yield sustained arrow-of-time in the autonomous accounted model.
To obtain nonzero steady-state entropy production, one needs a nontrivial affinity component (P6) in the lifted dynamics (e.g.\ a biased phase cycle) or an external schedule.
\end{corollary}

\begin{remark}[What the ``protocol trap'' actually warns about]
The theorem above is about \emph{true} path reversal asymmetry.
In practice, one can still obtain a spurious ``arrow'' by fitting a one-step Markov model to a non-Markov observed process (for example, by replacing the hidden-phase model by a stroboscopic kernel on $X$).
The cure is the clock audit: include the phase variable in the state, or work directly with path-space quantities that do not assume Markovianity at the observed level.
\end{remark}

\begin{remark}[Interpretation]
Defining irreversibility as path-space KL makes a decisive separation: coarse observation cannot create time-reversal asymmetry (no false positives), though it can hide it (false negatives).
The protocol trap is the concrete warning that an apparent arrow can be an artifact of hiding a clock/phase variable or fitting a one-step Markov model to a non-Markov observation.
In an autonomous lifted model, holonomy alone does not manufacture sustained directionality unless there is a genuine affinity (a P6-type bias) in the lifted dynamics.
\end{remark}

% Scope IDs: see docs/spec/theorem-inventory.md and paper-contract.md
\section{Generic extension and the finite forcing lemma}\label{sec:forcing}
\subsection{Languages as partitions and definability}

Assume \textbf{A\_FIN} and fix a finite microstate set $Z$.
A \emph{language} (or \emph{lens}) is a finite-valued map $f:Z\to X$ (assumption \textbf{A\_LENS}), inducing the partition
\[
\Pi_f := \{B_x\}_{x\in X}, \qquad B_x := f^{-1}(x).
\]
We interpret $\Pi_f$ as the collection of distinctions expressible at the current layer.

\begin{definition}[Predicates and definability]
A (Boolean) \emph{predicate} is a function $h:Z\to\{0,1\}$.
We say that $h$ is \emph{definable from the language $f$} (equivalently: $\Pi_f$-measurable) if it is constant on each block:
\[
\forall x\in X,\ \forall z,z'\in B_x,\quad h(z)=h(z').
\]
Equivalently, there exists $\tilde h:X\to\{0,1\}$ with $h=\tilde h\circ f$.
\end{definition}

Adjoining a new predicate $h$ refines the language by considering the joint lens
\[
(f,h): Z \to X\times\{0,1\}.
\]
This extension is \emph{trivial} precisely when $h$ was already definable from $f$.

\subsection{Counting lemma: definable predicates are rare}

Let $N:=|Z|$ and $K:=|X|$.

\begin{lemma}[Counting definable predicates]\label{lem:count-definable}
Assume \textbf{A\_FIN}+\textbf{A\_LENS}.
The number of predicates definable from $f:Z\to X$ is exactly $2^K$.
Consequently, if $h$ is sampled uniformly from $\{0,1\}^Z$, then
\[
\mathbb P(h \text{ is definable from } f)=\frac{2^K}{2^N}=2^{-(N-K)}.
\]
\end{lemma}

\begin{proof}
A definable predicate is determined by choosing one bit per block $B_x$ (the common value on that block), hence $2^K$ choices.
There are $2^N$ total predicates on $Z$, so the stated probability follows.
\end{proof}

\subsection{Finite forcing: generic extensions are non-definable}

The preceding lemma is the finite analogue of the forcing slogan: \emph{generic extensions add non-definable predicates}.
To emphasize the “genericity” aspect, we also record a simple sufficient condition for a predicate to be \emph{strongly novel}.

\begin{theorem}[Finite forcing lemma / physical forcing lemma]\label{thm:finite-forcing}
Assume \textbf{A\_FIN}+\textbf{A\_LENS}.
Let $f:Z\to X$ be a language with blocks $\Pi_f=\{B_x\}_{x\in X}$, and sample $h:Z\to\{0,1\}$ uniformly from $\{0,1\}^Z$ (i.i.d.\ fair bits across microstates).
Then:
\begin{enumerate}
\item \textbf{Non-definability is generic:} $\mathbb P(h \text{ is not definable from } f)=1-2^{-(N-K)}$.
\item \textbf{Block-splitting is typical when blocks are large:}
for each block $B_x$,
\[
\mathbb P(h \text{ is constant on } B_x)=2^{1-|B_x|},
\]
and by a union bound,
\[
\mathbb P(\exists x\in X\text{ such that }h\text{ is constant on }B_x)\ \le\ \sum_{x\in X} 2^{1-|B_x|}.
\]
In particular, if every block has size at least $m\ge 2$, then
\[
\mathbb P(h \text{ splits every block}) \ \ge\ 1 - K\,2^{1-m}.
\]
\end{enumerate}
\end{theorem}

\begin{proof}
Item (1) is Lemma~\ref{lem:count-definable}.
For item (2), on a block of size $|B_x|$, the only constant labelings are “all zeros” and “all ones”, each with probability $2^{-|B_x|}$, hence $2^{1-|B_x|}$ total.
The union bound gives the stated inequality, and the final display uses $|B_x|\ge m$ for all $x$.
\end{proof}

\begin{corollary}[Generic language extension is strict]\label{cor:strict-language-extension}
Assume \textbf{A\_FIN}+\textbf{A\_LENS}.
With probability $1-2^{-(N-K)}$, the refined lens $(f,h):Z\to X\times\{0,1\}$ is a \emph{strict} refinement of $f$ (i.e.\ it distinguishes at least one pair of microstates previously indistinguishable under $f$).
\end{corollary}

\begin{proof}
If the refinement were not strict, then $h$ would be constant on each block of $f$, i.e.\ definable from $f$.
The complement event has probability $1-2^{-(N-K)}$ by Theorem~\ref{thm:finite-forcing}(1).
\end{proof}

\begin{remark}[Why this is the anti-saturation move]
Closure (or packaging) at a fixed language saturates by idempotence: repeated closure steps do not create an infinite strict chain.
Corollary~\ref{cor:strict-language-extension} exhibits a mathematically clean \emph{extension step}: adjoining a generic predicate changes the language itself in a way that is overwhelmingly not reconstructible from the old distinctions.
This is the finite “forcing-style” mechanism we use to justify strict ladder climbing via language/closure change.
\end{remark}

\begin{remark}[Interpretation]
The forcing-style extension step is a finite counting phenomenon: when a language has hidden volume inside its blocks, the predicates definable from that language form an exponentially small subset of all predicates.
Therefore a random new predicate is overwhelmingly likely to be non-definable, giving a clean ``anti-saturation'' move: change the language, then re-close.
This is the finite proxy for the role that generic extensions play in infinitary logic.
\end{remark}

% Scope IDs: see docs/spec/theorem-inventory.md and paper-contract.md
\section{Primitives P1--P6 as closure-changing operations}\label{sec:primitives}
Primitives P1--P6 form a vocabulary of operations that change (i) the admissible transition
structure (support graph and cycle space), (ii) the available idempotent endomaps (packaging and
fixed points), and (iii) the affinity data encoded by the graph 1-form. These are not ratchets by
themselves; directionality is certified by path reversal asymmetry and/or non-exactness of the
log-ratio 1-form under the \textbf{A\_AUT}+\textbf{A\_REV}+\textbf{A\_ACC} regime.

\subsection{Definitions of P1--P6}

\begin{definition}[P1: operator rewrite]
P1 is a rewrite of the substrate operator: replace a Markov kernel $P$ by a new kernel $P'$ (or a
finite family $\{P^{(m)}\}$), thereby changing the endomap $\mu\mapsto \mu P^\tau$ and the induced
empirical endomap $E_{\tau,f}$. This is a change in the operator itself, not an external schedule.
\end{definition}

\begin{definition}[P2: gating / constraints]
P2 restricts the support graph by deleting edges (setting selected $P_{ij}=0$) and renormalizing
rows, or equivalently restricting to a subgraph on which the kernel lives. Cycle-rank claims are
interpreted on the undirected support graph justified by \textbf{A\_REV}.
\end{definition}

\begin{definition}[P3: autonomous protocol holonomy]
P3 is modeled in the autonomous lifted form on $Z:=X\times\Phi$. The phase $\phi\in\Phi$ evolves
by an internal kernel $S$, and conditioned on $\phi$ the microstate updates by $K_\phi$.
No external schedule is assumed in the main text; externally scheduled stroboscopic protocols are
non-autonomous and fall outside \textbf{A\_AUT}.
\end{definition}

\begin{definition}[P4: sectors / invariants]
P4 is a conserved sector label: the support decomposes into disconnected components (block
structure), or equivalently $P$ is block diagonal up to permutation, so evolution preserves a
sector coordinate.
\end{definition}

\begin{definition}[P5: packaging]
P5 is an idempotent endomap $e$ whose fixed points $\mathrm{Fix}(e)$ are the packaged objects at a
given layer. This is an endomap notion (Section~\ref{sec:idempotent-endo}), not an order-closure.
\end{definition}

\begin{definition}[P6: affinity / non-exact drive]
P6 is the presence of a nontrivial cohomology class of the antisymmetric log-ratio 1-form
$a_{ij}=\log(P_{ij}/P_{ji})$ on the bidirected support. P6 is ``on'' iff some cycle integral is
nonzero, i.e.\ $a$ is not exact.
\end{definition}

\subsection{Two load-bearing propositions}

\begin{theorem}[P2 gating shrinks cycle space (T-P2-01)]
Assume \textbf{A\_FIN}+\textbf{A\_REV}.
Let $G=(V,E)$ be an undirected support graph and let $G'=(V,E')$ be obtained by deleting edges.
Then
\[
\beta_1(G')\le \beta_1(G),
\qquad
\beta_1(G):=|E|-|V|+c(G),
\]
where $c(G)$ is the number of connected components.
\end{theorem}

\begin{proof}
Deleting a single edge reduces $|E|$ by one and increases $c(G)$ by at most one. Hence
$|E|-|V|+c(G)$ cannot increase. Iterating over all deleted edges yields the claim.
\end{proof}

\begin{theorem}[P1 can change cycle rank and spectral gap (T-P1-01)]
Assume \textbf{A\_FIN}+\textbf{A\_REV}.
P1 rewrites can increase or decrease cycle rank and can increase or decrease a metastability proxy
(the spectral gap of a reversible lazy random walk).
\end{theorem}

\begin{proof}[Proof sketch]
(Cycle rank.) Adding an edge within an existing connected component leaves $c(G)$ unchanged and
increases $|E|$ by one, so $\beta_1(G)$ increases by one. More general rewires can change $c(G)$
and thus increase or decrease $\beta_1$.

(Spectral gap.) By rewriting edge weights or topology, one can create or remove bottlenecks.
Highly connected graphs have large gaps, while nearly decomposable graphs have small gaps; thus
P1 can increase or decrease the gap by explicit constructions (see Appendix~C for examples).
\end{proof}

\begin{remark}[Evidence pointers]
Appendix~C records finite graph and kernel witnesses for the P1 and P2 behaviors described above,
including explicit cycle-rank changes and spectral-gap increases/decreases.
\end{remark}

\begin{remark}[Interpretation]
The primitives form a minimal control panel for what kinds of fixed points (objects) and obstructions (affinities) can exist.
P1/P2 change the transition structure (weights/topology), P5 extracts fixed points (packaging), and P6 is the genuine non-exactness that certifies sustained directional structure under autonomy and accounting.
Nothing in P1--P5 is teleological by itself: directionality is certified by path-space asymmetry and/or nonzero affinities, not by a coarse description that forgets variables.
\end{remark}

% Scope IDs: see docs/spec/theorem-inventory.md and paper-contract.md
\section{Examples}
\label{sec:examples}

This section gives a few finite examples that instantiate the preceding definitions and results.
They introduce no new machinery and are included purely for intuition and for checking that the
assumption tags match the intended use-cases.

\subsection{A biased three-cycle: a non-exact graph 1-form}
\label{subsec:ex:3cycle}

Let $Z=\{0,1,2\}$ and consider the Markov kernel $P$ defined by the parameters $(p,q,s)$ with
$p,q,s>0$ and $p+q+s=1$, where from each $i\in Z$,
\[
P(i,i)=s,\qquad P(i,i+1\bmod 3)=p,\qquad P(i,i-1\bmod 3)=q.
\]
This satisfies \textbf{A\_FIN}+\textbf{A\_REV}. The antisymmetric edge 1-form
$a_{ij}:=\log\!\big(P_{ij}/P_{ji}\big)$ is therefore defined on each bidirected edge.

Along the directed cycle $\gamma=(0\to 1\to 2\to 0)$, the cycle integral (affinity) is
\[
\oint_{\gamma} a
= \log\frac{P_{01}}{P_{10}}+\log\frac{P_{12}}{P_{21}}+\log\frac{P_{20}}{P_{02}}
= 3\log\frac{p}{q}.
\]
For instance with $(p,q,s)=(0.7,0.2,0.1)$ we get $|\oint_\gamma a|\approx 3.7583\neq 0$, hence the 1-form
is not exact and no global potential $\Phi$ can satisfy $a_{ij}=\Phi_j-\Phi_i$ on all edges (cf.\ Section~\ref{sec:acc}).

\subsection{Protocol trap: external schedule vs autonomous lifted model}
\label{subsec:ex:protocol-trap}

We illustrate the difference between an externally scheduled protocol and an autonomous lifted model.
Let $Z=\{0,1,2\}$ and let $K_0,K_1$ be two kernels that are each reversible w.r.t.\ the uniform distribution
(on $Z$), but do not commute (so the composition depends on order). One concrete instance used in our
reproducibility artifacts is:
\[
K_0=
\begin{pmatrix}
0.573333 & 0.393333 & 0.033333\\
0.393333 & 0.573333 & 0.033333\\
0.033333 & 0.033333 & 0.933333
\end{pmatrix},
\qquad
K_1=
\begin{pmatrix}
0.933333 & 0.033333 & 0.033333\\
0.033333 & 0.573333 & 0.393333\\
0.033333 & 0.393333 & 0.573333
\end{pmatrix},
\]
for which $\max_{i,j}|(K_0K_1-K_1K_0)_{ij}|\approx 1.296\times 10^{-1}$.

\emph{Externally scheduled (not autonomous):}
if one applies $K_0$ then $K_1$ in a fixed periodic order but hides the schedule variable, the
stroboscopic one-step kernel on $Z$ is $K:=K_1K_0$, which can exhibit positive steady-state entropy
production (positive time-reversal asymmetry in the one-step Markov description).

\emph{Autonomous lifted model (clock included in state):}
define an augmented state space $Z\times \Phi$ where $\Phi=\{0,1\}$ is a phase variable, and define a
time-homogeneous Markov chain that (i) updates $z\in Z$ using $K_\phi$ conditional on the current phase,
and (ii) updates $\phi$ using a Markov kernel $S$ on $\Phi$. If $S$ is unbiased/reversible (no phase
affinity) and both $K_0,K_1$ are reversible w.r.t.\ a common $\pi$, then the lifted model exhibits
vanishing steady-state entropy production (and hence vanishing path reversal asymmetry at stationarity),
whereas introducing a phase bias (nontrivial affinity in the $\Phi$-subsystem) yields strictly positive
entropy production. This is the clean mathematical form of the slogan “P3 needs P6” in autonomous models:
noncommuting actions alone do not create sustained directionality unless the protocol degree of freedom
itself is driven.

\subsection{Finite forcing count: definability is exponentially rare}
\label{subsec:ex:forcing-count}

Let $Z$ be a finite set with $|Z|=N$, and let a “language” be represented by a partition $\Pi$ of $Z$
into $K$ blocks (equivalently a lens $f:Z\to X$ with $|X|=K$). A predicate $h:Z\to\{0,1\}$ is definable
from $\Pi$ iff it is constant on each block of $\Pi$, so the number of definable predicates is exactly $2^K$,
while the total number of predicates is $2^N$. Therefore, if $h$ is sampled uniformly from $\{0,1\}^Z$,
\[
\mathbb P(h\text{ is definable from }\Pi)=\frac{2^K}{2^N}=2^{-(N-K)}.
\]
For example, when $(N,K)=(16,4)$, this gives $2^{-(16-4)}=2^{-12}\approx 2.44\times 10^{-4}$, i.e.\
a uniformly random new predicate is overwhelmingly likely to be non-definable from the old language.
This is the finite counting core behind the “generic extension” step in Section~\ref{sec:forcing}.

\section{Discussion and scope boundary}
\subsection{What the theory does and does not claim}\label{sec:discussion-claims}

The organizing loop in Section~\ref{sec:big-picture} separates three roles that are often conflated:
packaging (idempotence and fixed points), extension (non-definability), and directionality (path-space KL / affinities).
The paper’s contribution is to make these roles explicit and composable, with certificates that are stable under coarse observation (DPI) and robust to hidden-clock artifacts (protocol audit).

Equally important are the non-claims.
We do not claim that closure ladders arise automatically in arbitrary dynamics, nor that any particular choice of lens $f$ or timescale $\tau$ is preferred.
We also do not claim that protocol holonomy by itself yields sustained directionality under autonomy; on the contrary, the lifted-model statements formalize that sustained asymmetry requires a genuine affinity (a non-exact 1-form component) or an externally imposed schedule (which lies outside \textbf{A\_AUT}).
Finally, the finite forcing lemma is used only as a \emph{finite proxy} for generic extension: it justifies that strict language extension is cheap and typical when there is hidden volume, not that all forms of novelty reduce to random predicates.

\subsection*{Recap: what this paper establishes}

\paragraph{Summary of contributions.}
This paper isolates a small mathematical spine for ``closure ladder'' behavior in finite settings:
(i) order-theoretic closure ladders and the saturation obstruction from idempotence;
(ii) a unifying abstraction via idempotent endomaps and fixed points, allowing dynamics-induced operators $E_{\tau,f}$ to be treated without conflating them with order-closures;
(iii) a DTMC-friendly accounting regime in which affinities are cycle integrals of an antisymmetric log-ratio 1-form on the bidirected support graph;
(iv) an arrow-of-time diagnostic as path-space KL together with data processing, ensuring coarse-graining cannot create irreversibility;
and (v) a finite ``forcing'' lemma showing that generic predicate extensions are overwhelmingly non-definable from the old language, providing a clean anti-saturation move for strict ladder climbing.

\paragraph{How to interpret primitives without teleology.}
Primitives P1--P6 are a deliberately small \emph{complete vocabulary} of closure-changing operations in this framework: the theory is stated so that each primitive names a distinct kind of mathematical ingredient that can be present or absent, and together they span the mechanisms by which the ambient structure (and hence available fixed points/objects) can change.
P1 rewrites or replaces the operative kernel (changing the substrate operator itself); P2 gates feasibility by restricting support (edge deletion / renormalization); P3 introduces autonomous protocol holonomy by including a phase variable in the state (no external schedule in main-text claims); P4 introduces discrete sector structure/invariants (state-space decomposition); P5 is packaging as (approx) idempotent fixed-point extraction; and P6 is the non-exact affinity content (cycle obstruction) that certifies genuine driven directionality under autonomy/accounting.
None of P1--P5 is teleological on its own: directionality is certified by path-space asymmetry and/or nonzero affinities, and apparent arrows created by hiding variables or imposing an external schedule are separated from true irreversibility by the path-space definition and the protocol clock audit.

\paragraph{Limitations (by design).}
We stay in finite state spaces and emphasize graph-theoretic and measure-theoretic statements.
Computational artifacts in the repository serve only as reproducible sanity checks and are not used as proof premises.
We also avoid the continuous-time interpretation layer (stochastic thermodynamics identities, physical units, estimation/sample-complexity), which would require additional modeling choices outside the present math-only scope.

\appendix

\section{Appendix A: Reproducibility checklist}
This repository is intended to be mechanically checkable at three levels:
(i) a claim/dependency registry, (ii) a small Lean formalization core, and
(iii) a deterministic Python evidence harness.

\subsection*{A.1 Repository integrity checks (from repo root)}
Run the following commands from the repository root:
\begin{verbatim}
python3 scripts/check_deps_dag.py
python3 scripts/check_kb_pointers.py
python3 scripts/check_paper_contract.py
./check_lean.sh
\end{verbatim}

\subsection*{A.2 Python evidence harness (deterministic tests)}
The Python code lives under \texttt{python/} and is run in an isolated virtual environment.
From the repository root:
\begin{verbatim}
python3 -m venv python/.venv
python/.venv/bin/pip install -q pytest numpy
python/.venv/bin/pip install -e python
cd python && .venv/bin/python -m pytest -q
\end{verbatim}

\subsection*{A.3 Notes on tool availability}
A \LaTeX{} toolchain is not required to run the checks above.
If you wish to compile the PDF locally, install one of \texttt{latexmk}, \texttt{pdflatex}, or \texttt{tectonic} and run:
\begin{verbatim}
./build_paper.sh
\end{verbatim}

\section{Appendix B: Lean formalization map}
The Lean development provides a minimal mechanized backbone for the order-theoretic and
``thin packaging'' parts of the theory. It is intentionally small: no probability theory
or Markov chain results are mechanized in Lean for this paper.

\subsection*{B.1 What is mechanized}
The following items are implemented in Lean under \texttt{formal/}:

\begin{itemize}
  \item \textbf{Order-closure and closure ladders} (order-theoretic closure operators, fixed points, strict ladder relation, and iterate stabilization).
  \item \textbf{Thin packaging / reflection} for closure operators, expressed via the Galois insertion associated to a closure operator.
  \item \textbf{Idempotent endomaps} and fixed-point subtypes; bridge lemma from order-closures to idempotent endomaps.
\end{itemize}

\subsection*{B.2 File map and key declarations}
\begin{itemize}
  \item \texttt{formal/ClosureLadder/Basic.lean}:
    closure operators, closed subtype, strict ladder relation, and iterate stabilization.
    Key lemmas used in the manuscript include \texttt{ladder\*mono},
    \texttt{closure\*iterate\*succ}, and \texttt{closure\*iterate\*ge\*one}.
  \item \texttt{formal/ClosureLadder/Packaging.lean}:
    reflection-style lemmas for the closed subtype (thin-category packaging anchor).
  \item \texttt{formal/ClosureLadder/IdempotentEndo.lean}:
    the structure \texttt{IdempotentEndo}, its fixed-point subtype, and the bridge
    \texttt{ClosureOperator.toIdempotentEndo}.
  \item \texttt{formal/ClosureLadder.lean}:
    umbrella import module.
\end{itemize}

\subsection*{B.3 How to build}
From the repository root:
\begin{verbatim}
./check_lean.sh
\end{verbatim}
This script performs a cache fetch (when available) and then runs \texttt{lake build}.

\section{Appendix C: Python evidence map}
The repository contains a small deterministic Python harness (finite-state Markov kernels,
path laws, KL divergence, coarse-graining pushforwards, graph cycle utilities, and related
diagnostics). These computations are \emph{evidence and sanity checks} for the mathematical
definitions and theorem statements; they are not used as premises in any proof.

\subsection*{C.1 How to run}
See Appendix~A.2 for a one-shot test run via \texttt{pytest}. All tests are deterministic
(fixed seeds or no randomness).

\subsection*{C.2 Evidence by theme (tests and scripts)}
The following tests and scripts witness the main diagnostic phenomena discussed in the paper:

\paragraph{Path reversal asymmetry and data processing (DPI).}
Tests include randomized DPI checks, strict hiding examples, and nonstationary initial distributions.
Representative files:
\begin{verbatim}
python/tests/test_paths_kl_dpi.py
python/tests/test_paths_dpi_randomized.py
python/tests/test_paths_dpi_nonstationary_init.py
python/tests/test_paths_strict_hiding.py
python/scripts/dpi_sweep.py
\end{verbatim}

\paragraph{Protocol trap and the ``P3 needs P6'' correction under autonomy.}
Tests separate externally scheduled noncommuting kernels from autonomous phase-included models,
including m-phase protocols and the ``protocol clock audit''.
Representative files:
\begin{verbatim}
python/tests/test_protocol_trap_external_vs_autonomous.py
python/tests/test_protocol_clock_audit_phi_included.py
python/tests/test_protocol_mphase_random_scan.py
python/tests/test_protocol_stroboscopic.py
python/scripts/protocol_trap_demo.py
\end{verbatim}

\paragraph{Graph topology effects of P2 (edge deletion) and P1 (rewrites).}
Tests verify cycle-rank monotonicity under edge deletion, show cycle-rank increases under edge addition,
and demonstrate spectral-gap changes (including a decisive gap-decrease rewrite).
Representative files:
\begin{verbatim}
python/tests/test_graph_p2_cycle_rank_monotone.py
python/tests/test_graph_p1_edge_add_beta1_delta.py
python/tests/test_p1_spectral_gap_decrease.py
python/scripts/graph_p1_p2_demo.py
python/scripts/p1_gap_decrease_demo.py
\end{verbatim}

\paragraph{Dynamics-induced empirical endomap and idempotence defect.}
Tests implement the induced empirical endomap, the total-variation idempotence defect, and show that
refinement can help or hurt depending on leakage and timescale.
Representative files:
\begin{verbatim}
python/tests/test_empirical_closure_idempotence_defect_tv.py
python/tests/test_empirical_closure_refinement_reveals_objects.py
python/tests/test_empirical_closure_refinement_can_hurt.py
python/scripts/empirical_closure_demo.py
python/scripts/refinement_help_hurt_sweep.py
\end{verbatim}

\paragraph{Finite forcing / definability rarity.}
Tests include exact enumeration for small cases, Monte Carlo checks, and a contrast with fixed-weight sampling.
Representative files:
\begin{verbatim}
python/tests/test_forcing_exact_enumeration.py
python/tests/test_forcing_monte_carlo_matches_theory.py
python/tests/test_forcing_fixed_weight_differs.py
python/scripts/forcing_sweep.py
\end{verbatim}

\paragraph{ACC as graph 1-form exactness and cycle affinities.}
Tests compute the antisymmetric log-ratio 1-form, cycle integrals, and potential reconstruction when exact.
Representative files:
\begin{verbatim}
python/tests/test_acc_graph_one_form.py
\end{verbatim}

\bibliographystyle{alpha}
\bibliography{refs}

\end{document}
